\section{Communication protocols}
\label{sec:protocols}

Three contracts tie the tiers together: the BLE telemetry the collar advertises, the GATT
services for streaming and updates, and the over-the-air model protocol. Keeping these explicit
is what lets the collar and station firmware evolve independently.

\subsection{Telemetry advertisement}

In production mode the collar publishes its result in the manufacturer-specific field of its BLE
advertisement, so the station can read it passively without connecting. Version 2 of the contract
is shown in \cref{tab:telemetry}.

\begin{table}[ht]
  \centering
  \caption{Telemetry v2 manufacturer-data contract.}
  \label{tab:telemetry}
  \begin{tabularx}{\textwidth}{@{}llX@{}}
    \toprule
    \textbf{Byte} & \textbf{Field} & \textbf{Meaning} \\
    \midrule
    \texttt{mfg[2]} & mode       & \texttt{2} = real (on-device classification), \texttt{1} = simulated \\
    \texttt{mfg[3]} & class index & Predicted class, or \texttt{0xFF} for unknown/uncertain \\
    \texttt{mfg[4]} & confidence  & Confidence, \texttt{0}--\texttt{100} \\
    \bottomrule
  \end{tabularx}
\end{table}

Because the class is an index, not a name, the station and dashboard resolve it through the label
list for the running model version, keeping history correct across label-set changes
(\cref{sec:cloud-backend}).

\subsection{GATT services}

For training-mode streaming and for model updates the station connects to the collar. Streaming
carries audio and motion to the station for capture; the over-the-air service carries model data
to the collar. The OTA service uses a small family of \texttt{0x0b}\,\(\ast\)\(\ast\) UUIDs and a
status notification the collar uses to acknowledge progress and report success or failure.

\subsection{Over-the-air model protocol}

Model delivery is a simple, framed protocol shared by both firmwares through a common header.
A transfer is \texttt{BEGIN}, a stream of \texttt{DATA} chunks, then \texttt{END}; the
\texttt{BEGIN} frame declares the target slot, the model length, and a CRC so the collar can
reject a corrupt transfer before committing it.

\begin{lstlisting}[language=C,caption={Shared OTA framing (\texttt{firmware/common/meow\_ota.h}, illustrative).},label={lst:ota}]
enum meow_ota_op   { MEOW_OTA_BEGIN = 1, MEOW_OTA_DATA = 2, MEOW_OTA_END = 3 };
enum meow_ota_slot { MEOW_SLOT_IMU  = 0, MEOW_SLOT_AUDIO = 1 };

struct meow_ota_begin {
    uint8_t  op;      // MEOW_OTA_BEGIN
    uint8_t  slot;    // MEOW_SLOT_IMU | MEOW_SLOT_AUDIO
    uint32_t len;     // model size in bytes
    uint32_t crc32;   // CRC-32/IEEE over the model blob
} __packed;
\end{lstlisting}

On the collar, \texttt{END} verifies the CRC, writes the model header, and sets a pending flag;
it does \emph{not} swap the live model from inside the BLE callback. The actual swap is performed
later by a service call from the main loop, on the main thread's larger stack, because committing
a model is stack-heavy and doing it in the radio callback crash-looped the device. This split,
together with the raised main stack (\cref{sec:firmware-collar}), is what made OTA reliable on
hardware.
